\section{Применение к данным}
\label{sec:results}

\subsection{Прямая задача: data1 и опубликованное решение}
\draft{Для data1 доступно опубликованное эталонное решение. Сравнение
выполнено на исходной геометрии XLSX по значениям на кровле, подошве
и средним по слою. На рис.~\ref{fig:data1-temperature}
и~\ref{fig:data1-ro} показаны современные поля и знаковые разности.
Используется относительная табличная норма без пространственных весов:}
\begin{equation}
 \varepsilon_{\mathrm{tab}}(u)=100
 \frac{\left(\sum_j(u_j^{\mathrm{calc}}-u_j^{\mathrm{ref}})^2\right)^{1/2}}
      {\left(\sum_j(u_j^{\mathrm{ref}})^2\right)^{1/2}}\quad[\%].
\end{equation}
\reportfigure{12_data1_temperature.png}
  {Data1, современная температура: эталон, расчёт и расчёт минус эталон.
  Поля восстановлены по значениям на кровле и подошве каждого слоя.}
  {fig:data1-temperature}
\reportfigure{13_data1_reflectance.png}
  {Data1, современная зрелость Ro: эталон, расчёт и их знаковая разность.}
  {fig:data1-ro}

\draft{По средним значениям слоёв ошибки T/Ro за все 18 возрастов равны
2,803/1,859\%; для современной конфигурации --- 3,598/2,385\%.
Рисунок~\ref{fig:profiles} показывает вертикальные профили у области
наибольшего расхождения и вдали от неё.}
\reportfigure{14_data1_vertical_profiles.png}
  {Data1: вертикальные профили T и Ro при $x=282$ и $1400$ м;
  синие линии --- расчёт, оранжевые --- эталон.}
  {fig:profiles}
\FloatBarrier

\subsection{Прямая задача: data2 и упрощение геометрии}
\label{sec:data2}
\draft{На рис.~\ref{fig:data2} показан второй прямой набор,
для которого эталон не опубликован. На нём исследуется разделение
исходной геометрии и расчётной сетки: упрощение полилиний RDP сокращает
число колонок, после расчёта решение возвращается на исходные координаты.}
\reportfigure{15_data2_fields.png}
  {Data2: современные поля T и Ro с границами слоёв.}
  {fig:data2}
\draft{Рисунок~\ref{fig:rdp} показывает ускорение и связанное с ним изменение
решения: при допуске 5 м число ячеек уменьшается с 120335 до 3943,
время --- с 205,4 до 5,0 с. Относительные табличные отличия средних T/Ro
равны 1,31/1,12\%, но локальное отличие Ro достигает 6,91\%.
Это одиночные замеры и сравнение с расчётом без упрощения;
геометрия и сетка меняются одновременно.}
\reportfigure{24_geometry_simplification_tradeoff.png}
  {Data2: влияние допуска упрощения границ на размер сетки,
  время и отличия современных T и Ro от расчёта без упрощения.}
  {fig:rdp}
\FloatBarrier

\subsection{Обратная задача}
\draft{Начальный и найденный потоки показаны на рис.~\ref{fig:inverse-q}.
Изменение потока на последней эпохе требует дальнейшей проверки весов
и размерности параметризации; найденная история не объявляется единственной.}
\reportfigure{16_inverse_basal_flux.png}
  {Начальное приближение и найденный нижний поток в 17 узлах
  с линейной интерполяцией между ними.}
  {fig:inverse-q}

\draft{Подбор выполнен на сетке $200\times100$ м: 10 итераций,
44 прямых вычисления, $J=265{,}883$; остановка по лимиту итераций.
Рисунок~\ref{fig:optimization} показывает лучшее значение функционала
и стоимость начальной серии из 18 вычислений. Контроль найденного потока
на сетке $200\times50$ м даёт $J_T=205{,}898$, $J_{\Ro}=55{,}608$.}
\reportfigure{17_inverse_objective.png}
  {История функционала по числу прямых вычислений; выделена начальная серия оптимизатора.}
  {fig:optimization}

\draft{Итоговые профили по пяти скважинам приведены на
рис.~\ref{fig:wells-temperature} и~\ref{fig:wells-ro}.
В допуске 456/500 температур и 39/50 значений Ro.
Среднеквадратичные относительные невязки --- 1,283\% и 2,109\%,
максимальные --- 6,743\% и 6,907\% соответственно.}
\reportfigure{18_inverse_wells_temperature.png}
  {Температура в скважинах A--E при найденном потоке:
  расчёт, измерения и интервалы погрешности 2\%.}
  {fig:wells-temperature}
\reportfigure{19_inverse_wells_reflectance.png}
  {Зрелость Ro в тех же скважинах: расчёт, измерения и интервалы погрешности 2\%.}
  {fig:wells-ro}

\draft{Рисунок~\ref{fig:inverse-fields} показывает современные T и Ro
во всём разрезе при том же потоке, включая участки вне скважин.
Для этих полей эталон неизвестен.}
\reportfigure{20_inverse_fields.png}
  {Полные современные поля T и Ro обратного набора при найденном потоке.}
  {fig:inverse-fields}
\draft{При том же найденном потоке отдельно изменяются $\Delta z$, $\Delta t$
и $\varepsilon$ (рис.~\ref{fig:sensitivity}). Переход с 1 на 0,25 млн лет
меняет T в точках наблюдений не более чем на $0{,}0013\,{}^\circ$C,
Ro --- на 0,120\%. Это чувствительность рабочего расчёта к дискретизации.}
\reportfigure{11_inverse_discretization_sensitivity.png}
  {Изменения T и Ro в точках наблюдений при изменении одного
  параметра дискретизации относительно основного расчёта.}
  {fig:sensitivity}
\FloatBarrier

\subsection{Время расчёта}
\label{sec:runtime}
\draft{Измерения: Intel Core i7-12700H, один вычислительный поток,
GCC 15.2, оптимизированная сборка, PETSc 3.24.4;
$\Delta x=200$ м, $\Delta z=50$ м, $\Delta t=1$ млн лет,
$\varepsilon=0{,}05$ м, допуск упрощения 5 м.
Относительные допуски линейного и нелинейного решений --- $10^{-8}$.
Рисунок~\ref{fig:timing} показывает распределение времени по этапам.
Подготовка входа и экспорт XLSX в Python в эти времена не входят.}
\reportfigure{21_forward_stage_timings.png}
  {Доли времени прямого расчёта на трёх наборах:
  стационарный старт, сетка, перенос, сборка, линейное решение, Ro и вывод.}
  {fig:timing}
\draft{На рис.~\ref{fig:solvers} GMRES(100)/ILU(2) сопоставлен с UMFPACK~\cite{suitesparse}.
По медианам трёх запусков ускорение линейного решения составляет 4,05;
3,67 и 1,44 раза для data1, data2 и обратного набора соответственно;
всего прямого расчёта --- 2,55; 2,32 и 1,31 раза.
Различия T не превышают $3\cdot10^{-11}\,{}^\circ$C.}
\reportfigure{22_linear_solver_comparison.png}
  {Время линейного решения с подготовкой и полного прямого расчёта:
  GMRES(100)/ILU(2) и UMFPACK, медианы и диапазоны трёх запусков.}
  {fig:solvers}
\draft{Изолированное вычисление интегралов EASY\%Ro с таблицей
быстрее прямой аналитики через интегральную экспоненту в 6,52 раза
(рис.~\ref{fig:easy-speed}). Однократное построение таблицы занимает около
0,25 мс. Это ускорение вычисления интегралов, а не всего прямого расчёта.}
\reportfigure{23_easy_ro_table_benchmark.png}
  {Стоимость вычисления интегралов через аналитическую первообразную
  и таблицу; справа --- максимальная относительная ошибка Ro
  на четырёх историях относительно независимой квадратуры.}
  {fig:easy-speed}
